\chapter{Recursos e ambiente de desenovolvimento}\label{sec:5}

Nesta seção se apresenta todas as tecnologias que serão utilizadas durante o desenvolvimento do projeto. 

\section{Linguagens de Programação}

O projeto utilizará duas linguagens de programação complementares. No lado 
servidor (\textit{back-end}), será empregado o Java, linguagem 
orientada a objetos amplamente consolidada no mercado corporativo \cite{Oracle2025}. No lado 
cliente (\textit{front-end}), será utilizado o JavaScript, linguagem 
interpretada pelos navegadores web e amplamente adotada \cite{Guo2025}.

\section{Banco de Dados}

O sistema de gerenciamento de banco de dados adotado será o 
PostgreSQL \cite{PostgreSQLGlobalDevelopmentGroup2025}. A escolha se justifica pela maturidade da tecnologia e pela sua 
ampla compatibilidade com os \textit{frameworks} utilizados no projeto.

\section{\textit{Frameworks} de Desenvolvimento}

\subsection{\textit{Front-end}: React}

Para o desenvolvimento da interface do usuário, será utilizado o 
React, uma biblioteca voltada para a construção de interfaces reativas e componentizadas \cite{ReactTeam2025}. 

\subsection{\textit{Back-end}: Spring Boot}

Para o desenvolvimento da camada de negócios e exposição de APIs REST, será 
utilizado o Spring Boot, \textit{framework} baseado no ecossistema 
Spring para a plataforma Java \cite{Webb2025}.